\section{引言：从 Long CoT 到 Interleaved Thinking}

自 2024 年 9 月 OpenAI 推出 o1 (Reasoning Model) 以来，以及 2025 年 1 月 DeepSeek R1 的发布，``Thinking Model'' 成为大模型领域的重要范式。然而，传统的 Long Chain-of-Thought (Long CoT) 存在明显不足。

K2 Thinking 提出了一种全新的范式------\textbf{Interleaved Thinking}（交错推理/交错思考），面向 Agentic 场景，通过强化学习训练模型将思考与回答交替进行。

\begin{importantbox}{传统 Long CoT 的问题}
\begin{itemize}
  \item 在最终 answer 之前有一个漫长的 long-CoT 过程
  \item 首字延迟（Time to First Token, TTFT）极高，用户体验差
  \item 难以定义细粒度的 reward，过程监督困难
\end{itemize}
\end{importantbox}

